%% ELEMENTY

\usepackage[x11names]{xcolor}

\usepackage{microtype}

\usepackage[all]{nowidow}

\usepackage{graphicx}

\usepackage[hidelinks, pdflang=pl]{hyperref}
\hypersetup{
    colorlinks=false,
    linktoc=all,
} % linki w spisie treści

%% FONTY

\usepackage{fontspec}% font font font

\setmainfont{Poltawski Nowy}% font
\newfontfamily{\naut}{Poltawski Nowy}[Numbers=OldStyle]% jeśli potrzeba cyfr nautycznych
\newfontfamily{\linin}{Poltawski Nowy}[Numbers=OldStyleOff]
%% Inne proponowane:
% Courier Prime – Typewriter
% Coelacanth – Humanist antiqua
% Alegreya – Humanist
% Liberation Serif – Szeryfowa, Timesopodobna
% Noto Serif – Timesopodobna, obsługuje SC
% Poltawski Nowy – Polski klasyk

%% SKŁAD

\usepackage{sectsty}% zdefiniowanie wyglądu sekcji
% \allsectionsfont{\centering}%

\usepackage{geometry}
%% Wybierz format
\geometry{a4paper, bottom=30mm, top=30mm}% wszystkie wartości jeszcze mogą sie zmienić

\newcommand{\tytułowa}{\fontsize{40}{60}\selectfont}
\newcommand{\dziełowa}{\fontsize{12}{18}\selectfont}
\newcommand{\indeksowa}{\fontsize{10}{13}\selectfont}

\usepackage{epigraph}

\setcounter{secnumdepth}{-2}

\makeatletter
\renewcommand\tableofcontents{%
    \if@twocolumn
      \@restonecoltrue\onecolumn
    \else
      \@restonecolfalse
    \fi
    \section*{\contentsname%TUTAJ ZMIENIAĆ RANGĘ SPISU TREŚCI
        \@mkboth{%
           \MakeUppercase\contentsname}{\MakeUppercase\contentsname}}%
    \indeksowa\@starttoc{toc}%
    \if@restonecol\twocolumn\fi
}\dziełowa
\makeatother

\addto\captionspolish{\renewcommand{\contentsname}{Spis treści}}% zmiana nazwy spisu treści

\usepackage{fancyhdr}

\renewcommand{\headrulewidth}{0pt}
\fancyhead[L,C,R]{}
\fancyfoot[L,C,R]{}
\fancyfoot[C]{\thepage}

\usepackage{marginnote}% żeby móc zmieniać wysokość marginpar-ów

\usepackage{bbding}
%% WYMAGANE: graphicx, xparse, xstring, xfp
%% ZMIENNE:
\newcommand{\marginaliaraisebox}{0.2}%o ile podnieść manikułę;
% Coelacanth: 0.2
\newcommand{\marginaliawysokość}{0.1}%modyfikacja domyślnej wartości 

%% GLIF/BRAK GLIFU
% \PackageError{marginalia}{Nie wskazano czy font ma glif manikuły}{}
% \newcommand{\maniculeL}{\raisebox{\marginaliaraisebox cm}{\scalebox{0.25}{\Huge☜}}}\newcommand{\maniculeR}{\raisebox{\marginaliaraisebox cm}{\scalebox{0.25}{\Huge☞}}}%% DLA FONTÓW Z GLIFEM
\newcommand{\maniculeL}{\HandLeft}\newcommand{\maniculeR}{\HandRight}%% DLA FONTÓW BEZ GLIFU

\NewDocumentCommand{\Mar}{s d<> O{0} m O{} D<>{\marginaliawysokość}}{% #1: Boolean; #2: <P(rawy)/L(ewy)>; #3: [X° (obrót palca)]; #4: {<tytuł>/<treść>}; #5: [<treść>]; #6: <<zmiana wysokości>>
    \IfNoValueTF{#2}{%UNIWERSALNY MARGINPAR
        \ifodd\value{page}%PRAWY MARGINES
        \marginnote{{\kolor\scalebox{0.9}{\rotatebox[origin=c]{\inteval{-1 * #3}}{\maniculeL}}\IfBooleanT{#1}{\raisebox{0.06cm}{ #4}}}\fontsize{9}{9}\selectfont\\\vspace{1pt}\IfBooleanTF{#1}{#5}{#4}}[#6cm]%
        \else%LEWY MARGINES
        \marginnote{{\kolor\IfBooleanT{#1}{\raisebox{0.06cm}{#4 }}\scalebox{0.9}{\rotatebox[origin=c]{#3}{\maniculeR}}}\fontsize{9}{9}\selectfont\\\vspace{1pt}\IfBooleanTF{#1}{#5}{#4}}[#6cm]%
        \fi%
    }{%WYBRANY MARGINPAR
        \IfStrEqCase{#2}{%
            {P}{%PRAWY MARGINES
                \ifodd\value{page}%
                \else%
                \reversemarginpar%
                \fi%
                \marginnote{{\kolor\scalebox{0.9}{\rotatebox[origin=c]{\inteval{-1 * #3}}{\maniculeL}}\IfBooleanT{#1}{\raisebox{0.06cm}{ #4}}}\fontsize{9}{9}\selectfont\\\vspace{1pt}\IfBooleanTF{#1}{#5}{#4}}[#6cm]%
            }%
            {L}{%LEWY MARGINES
                \ifodd\value{page}%
                \reversemarginpar%
                \else%
                \fi%
                \marginnote{{\kolor\IfBooleanT{#1}{\raisebox{0.06cm}{#4 }}\scalebox{0.9}{\rotatebox[origin=c]{#3}{\maniculeR}}}\fontsize{9}{9}\selectfont\\\vspace{1pt}\IfBooleanTF{#1}{#5}{#4}}[#6cm]%
            }%
        }%
    }%
    \normalmarginpar%
\ }

%\Mar<L/P>[X°]{<treść>}<wysokość w cm>
%\Mar*<L/P>[X°]{<tytuł>}[<treść>]<wysokość w cm>

% Dla artykułu może być konieczne zakomentowanie odpowiednich fragmentów (\ifodd\value{page}) w części "WYBRANY MARGINES"


%% ZMIANY

\usepackage{chngcntr}

%% Tak jak w skrypcie z zagadnieniami z admina
%% Lepiej wygląda z tytułami wyrównanymi do lewej!

\newcounter{zagsectionc}
\setcounter{zagsectionc}{0}

\newcounter{zagc}
\setcounter{zagc}{0}

%% WYBIERZ SPOSÓB NUMERACJI:
% \PackageError{Komendy}{Wybierz sposób numeracji}{}

% %% 1; 2; 3; 4; 5; 6; 

\NewDocumentCommand{\zagsection}{O{#2} m}{%
    \refstepcounter{zagsectionc}% 
    \oldsection*{#2}% bez numeru w nazwie
    \phantomsection\label{zagsec:\arabic{zagsectionc}}% label z numerem sekcji (zagsec:n)
    \addcontentsline{toc}{section}{\kolor#1}% dodanie do spisu treści
}

\NewDocumentCommand{\zag}{O{#2} m O{}}{%
\refstepcounter{zagc}% 
\oldsubsection*{{\reversemarginpar\marginnote{\arabic{zagc}.\\#3 }}#2}% bez numeru w nazwie
\phantomsection\label{zag:\arabic{zagc}}% label z numerem subsekcji (zag:n)
\addcontentsline{toc}{subsection}{\fontsize{10}{10}\selectfont\arabic{zagc}. #1}% dodanie do spisu treści
% Opcjonalnie, jeśli chcę zmienić nagłówek:
\markboth{\textup{Zagadnienie \thezagc}}{\textup{Część \Roman{zagsectionc}}}%
}% Stała numeracja zagadnień
% 
%% 1.1; 1.2; 1.3; 2.1; 2.2; 3.1

\NewDocumentCommand{\zagsection}{O{#2} m}{%
    \refstepcounter{zagsectionc}% 
    \setcounter{zagc}{0}% reset licznika subsekcji (1.1, 1.2 -> 2.1, 2.2 a nie 2.3)
    \oldsection*{#2}% bez numeru w nazwie
    \phantomsection\label{zagsec:\arabic{zagsectionc}}% label z numerem sekcji (zagsec:n)
    \addcontentsline{toc}{section}{#1}% dodanie do spisu treści
}

\NewDocumentCommand{\zag}{O{#2} m O{}}{%
\refstepcounter{zagc}% 
\oldsubsection*{{\reversemarginpar\marginnote{\arabic{zagc}.\\#3 }}#2}% bez numeru w nazwie
\phantomsection\label{zag:\arabic{zagsectionc}.\arabic{zagc}}% label z numerem subsekcji (zag:n.m)
\addcontentsline{toc}{subsection}{\fontsize{10}{10}\selectfont\arabic{zagc}. #1}% dodanie do spisu treści
% Opcjonalnie, jeśli chcę zmienić nagłówek:
\markboth{\textup{Zagadnienie \thezagc}}{\textup{Część \Roman{zagsectionc}}}%
}% Numeracja zagadnień resetowana co sekcję

%% DOKUMENTACJA:

%\zagsection[<toctytuł>]{<tytuł>}

%\zag[<toctytuł>]{<tytuł>}[<podpis>]
%% Tytuły zrównane do środka, § – subsekcja, dział – sekcja
\newcounter{umsectionc}
\newcounter{umparc}
\setcounter{umsectionc}{0}
\setcounter{umparc}{0}

\NewDocumentCommand{\umsection}{O{#2} m}{%
    \refstepcounter{umsectionc}%
    \oldsection*{\Roman{umsectionc}. #2}%
    \phantomsection\label{umsec:\roman{umsectionc}}%
    \addcontentsline{toc}{section}{\Roman{umsectionc}. #1}%
    \markboth{\textup{\@title}}{\textup{\Roman{umsectionc}. #1}}
}

\NewDocumentCommand{\umpar}{O{#2} m}{%
    \refstepcounter{umparc}%
    \oldsubsection*{§ \arabic{umparc}. #2}%
    \phantomsection\label{umpar:\arabic{umparc}}%
    \addcontentsline{toc}{subsection}{§ \arabic{umparc}. #1}%
}

\usepackage{multirow}

\newenvironment{newindent}[2]{\leftskip#1
}{\leftskip0cm}

\NewDocumentCommand{\art}{m o}{\noindent\textbf{Art.\ #1.%
\IfValueT{#2}{\ [#2]}}}

\NewDocumentCommand{\ust}{s m}{\IfBooleanTF{#1}{\noindent}{\\}\textbf{#2.}}

\newcommand{\pkt}[2]{\begin{newindent}{.25cm}{#1}

    \noindent\settowidth{\hangindent}{#1\ }\textbf{#1} #2

\end{newindent}}

\newcommand{\lit}[2]{\begin{newindent}{.75cm}{#1}

    \noindent\settowidth{\hangindent}{#1\ }\textbf{#1} #2

\end{newindent}}

\newcommand{\tiret}[2]{\begin{newindent}{1.25cm}{#1}

    \noindent\settowidth{\hangindent}{#1\ }\textbf{#1} #2

\end{newindent}}

%       SKŁADNIA
%
%       \art{numer}[nazwa]
%
%       \ust*{numer} noindent przy 1 
%       \ust{numer} nowa linia przy 1
%
%       \pkt{numer}{cała treść}
%
%       \lit{numer}{cała treść}
%
%       \tiret{numer}{cała treść}
%
%   zawsze przed \pkt \lit i \tiret dawaj pustą linijkę!
%       ^(.+\n)(\\pkt)
%       $1\n$2